\scanpage{91}

On peut supposer déjà $F$ et $F'$ isomorphes, \deleted{\unsure{si}} s'il est vrai que les deux foncteurs fibre sont isomorphes \unsure{fppf|fpqc}. --- Il semble que ça marche du moins \deleted{en car $0$} \uline{si $k$ est de car $0$}, en supposant tjrs $\mathcal M$ semi-simple ($Q$ apparaît comme $\varprojlim$ de schémas affines et plats sur $S$, donc est affine et plat).

% slanted
\margin{Construction en termes de $G'$ \deleted{\unsure{et|\({}^{\mathrm{ad}}\)}} $G_m\xrightarrow{i}G'$ (on oublie $G$)}

Déterminons le groupe $H'$, et pour ceci choisissons un \uline{générateur} $M$, alors $G'\subset\underline{\mathrm{Aut}}_{\underline{\mathcal O}_S}(F'(M))$, et $H'$ est le sous-groupe de $G'$ formé des $g'$ tels que, $\forall p\in\mathbb Z$,

\[
g'x_{p}^{\deleted{(p)}}\equiv x_{p}^{\deleted{\unsure{(p)}}}
\deleted{\unsure{+\,g_p x_{p'}^{(p')}}}
\inserted{\bigl(F'^{(p+1)}(M)\bigr)}
\deleted{x_{p}^{\unsure{(p)}}\in}
\]

Les transformations infinitésimales de $H'$ s'obtiennent ainsi : comme $G_m$ opère sur $\mathfrak g'$ grâce à $i:G_m\longrightarrow G'$, $\mathfrak g'$ se décompose suivant les caractères de $G_m$, et on a

\deleted{$\unsure{\mathfrak g}$}

\[
\boxed{\mathfrak h'=\sum_{n\geq1}\mathfrak g'_n}
\]

% slanted
\margin{\deleted{A examiner : si $G'$ \unsure{est lisse}, $H'$ \unsure{l'est-il|l'est} également ??}}

[N.B. Le sous-groupe qui invarie la filtration sur $\unsure{D'|S'}$ correspond à l'Algèbre de Lie $\mathfrak b'=\sum_{n\geq0}\mathfrak g'_n$ --- c'est le produit semi-direct du centralisateur de $i$ dans $\mathfrak g$, et de $\mathfrak h'\ldots$].

% slanted
\margin{N.B. $H'$ est \uline{tjrs} \uline{\unsure{connexe|ouvert|connexes}} (\uline{étant} \uline{unipotent} !)}

\deleted{Le groupe} L'algèbre $\mathfrak h'$ est nilpotente, et correspond d'ailleurs à un groupe \inserted{connexe} nilpotent, tandis que $\mathfrak b'$ est parabolique ; le centralisateur de $i$ \deleted{est} est un sous-groupe de Levi.

Pb : étudier les relations entre ``ss-groupes à $1$ paramètre'' et sous-groupes paraboliques avec déc de Lévi.
% archive sheet number: 90
